\documentclass[11pt,a4paper]{article}

\usepackage[utf8]{inputenc}
\usepackage[T1]{fontenc}
\usepackage[spanish]{babel}
\usepackage[margin=2.5cm]{geometry}
\usepackage{booktabs}
\usepackage{tabularx}
\usepackage{listings}
\usepackage{xcolor}
\usepackage{amsmath}
\usepackage{enumitem}
\usepackage{hyperref}

\hypersetup{
    colorlinks=true,
    linkcolor=black,
    urlcolor=blue,
    citecolor=black,
    pdftitle={Documento Tecnico - Libre Mercado Distribuido}
}

\lstset{
    basicstyle=\ttfamily\small,
    breaklines=true,
    columns=fullflexible,
    frame=single,
    backgroundcolor=\color{gray!8},
    xleftmargin=0pt,
    xrightmargin=0pt
}

\newcommand{\code}[1]{\texttt{#1}}

\title{\textbf{Documento Técnico --- Libre Mercado Distribuido}}
\author{Sistemas Distribuidos --- 2026 \\ % TODO: agregar nombre(s) del/los autor(es)
\textit{Nombre Apellido}}
\date{\today}

\begin{document}

\maketitle
\tableofcontents
\newpage

% ============================================================
\section{Arquitectura del Sistema}
% ============================================================

\subsection{Visión General}

Libre Mercado implementa una arquitectura distribuida de \textbf{3 nodos de base de datos independientes}, uno por sucursal, coordinados por una única capa de aplicación PHP. Cada nodo es una instancia MariaDB ejecutándose en su propio contenedor Docker, aislada de las demás.

\begin{lstlisting}
                    +----------------------+
                    |      Cliente Web     |
                    |  (Navegador / AJAX)  |
                    +-----------+----------+
                                | HTTP
                    +-----------v----------+
                    |     PHP + Apache     |
                    |  (Nodo Coordinador)  |
                    |   app:8080 (Docker)  |
                    +--+--------+--------+-+
                       |        |        |   PDO (red tienda_net)
              +--------v--+ +---v-----+ +-v---------+
              |  db_suc1  | | db_suc2 | |  db_suc3  |
              |  MariaDB  | | MariaDB | |  MariaDB  |
              | Suc.Norte | |Suc.Centro| | Suc. Sur |
              | BD Local  | | BD Local | | BD Local |
              +-----------+ +---------+ +-----------+
\end{lstlisting}

\subsection{Componentes}

\begin{table}[h]
\centering
\begin{tabularx}{\textwidth}{@{}l l X@{}}
\toprule
\textbf{Componente} & \textbf{Tecnología} & \textbf{Función} \\
\midrule
Frontend & HTML5 + JavaScript (AJAX) & Interfaz de tienda y administración \\
Coordinador & PHP 8.2 + Apache & Lógica de negocio, coordinación de transacciones \\
Nodo 1 & MariaDB 10.4 (\code{db\_suc1}) & Base de datos Sucursal Norte \\
Nodo 2 & MariaDB 10.4 (\code{db\_suc2}) & Base de datos Sucursal Centro \\
Nodo 3 & MariaDB 10.4 (\code{db\_suc3}) & Base de datos Sucursal Sur \\
Red & Docker bridge (\code{tienda\_net}) & Comunicación entre contenedores \\
\bottomrule
\end{tabularx}
\end{table}

\subsection{Archivos Clave}

\begin{table}[h]
\centering
\begin{tabularx}{\textwidth}{@{}l X@{}}
\toprule
\textbf{Archivo} & \textbf{Rol} \\
\midrule
\code{ConexionNodos.php} & Pool de conexiones PDO a los 3 nodos; detecta fallas simuladas y reales \\
\code{MultiNodo.php} & Ejecutor de operaciones atómicas sobre los 3 nodos (2PC para catálogo) \\
\code{api/procesar\_venta.php} & Llama a \code{sp\_realizar\_compra()} en el nodo de la sucursal \\
\code{api/traspaso\_stock.php} & Two-Phase Commit entre dos nodos para traspasos de inventario \\
\code{api/nodo\_estado.php} & Activa/desactiva fallas simuladas (Requisito 4) \\
\code{api/reconstruir\_stock.php} & Llama a \code{sp\_reconstruir\_stock()} tras recuperar un nodo \\
\code{api/ajustar\_stock.php} & Llama a \code{sp\_actualizar\_stock()} para ajustes manuales de inventario \\
\code{sql/suc1.sql}, \code{suc2.sql}, \code{suc3.sql} & Schema, datos iniciales y procedimientos almacenados \\
\bottomrule
\end{tabularx}
\end{table}

\newpage
% ============================================================
\section{Modelo Distribuido}
% ============================================================

\subsection{Estrategia de Distribución de Datos}

Se usa una combinación de \textbf{particionamiento} y \textbf{replicación} según el tipo de dato:

\textbf{Datos particionados (locales por nodo):}
\begin{itemize}[nosep]
    \item \code{stock} --- cada nodo tiene su propio inventario físico
    \item \code{ventas} y \code{detalle\_ventas} --- cada sucursal registra solo sus ventas
    \item \code{compras} y \code{detalle\_compras} --- cada sucursal gestiona su reabastecimiento
\end{itemize}

\textbf{Datos replicados (idénticos en los 3 nodos):}
\begin{itemize}[nosep]
    \item \code{productos} --- catálogo de artículos
    \item \code{usuarios} y \code{clientes} --- base de clientes
    \item \code{proveedores} --- catálogo de proveedores
    \item \code{sucursales} --- lista maestra de sucursales
\end{itemize}

La replicación del catálogo se mantiene mediante operaciones atómicas multi-nodo coordinadas por \code{MultiNodo.php}: toda creación, edición o baja en catálogo se aplica a los 3 nodos dentro de una misma transacción distribuida; si algún nodo falla, se hace rollback completo en todos.

\subsection{Flujo de una Compra Distribuida}

\begin{lstlisting}
Cliente -> selecciona producto -> agrega al carrito -> confirma compra
                                                          |
                             +----------------------------v------------------+
                             |          PHP: procesar_venta.php              |
                             |                                               |
                             |  CALL sp_realizar_compra(suc, cli, prod, cant)|
                             |                                               |
                             |  +-----------------------------------------+  |
                             |  |  MariaDB - Stored Procedure             |  |
                             |  |  START TRANSACTION                      |  |
                             |  |  SELECT cantidad ... FOR UPDATE         |  |
                             |  |  IF stock < cantidad -> ROLLBACK        |  |
                             |  |  INSERT INTO ventas ...                 |  |
                             |  |  INSERT INTO detalle_ventas ...         |  |
                             |  |  UPDATE stock SET cantidad - cant       |  |
                             |  |  COMMIT                                 |  |
                             |  +-----------------------------------------+  |
                             +-----------------------------------------------+
                                        |                    |
                                  [Exito: 200]         [Error: 409/503]
                                 Stock actualizado      Rollback automatico
\end{lstlisting}

\subsection{Flujo de un Traspaso (Two-Phase Commit)}

\begin{lstlisting}
Admin -> formulario traspaso -> envia a traspaso_stock.php
                                        |
                    +-------------------v------------------------+
                    |         Two-Phase Commit (2PC)              |
                    |                                              |
                    |  FASE 1 - PREPARE                            |
                    |  - beginTransaction() en nodo ORIGEN        |
                    |      UPDATE stock SET cantidad - X           |
                    |  - beginTransaction() en nodo DESTINO       |
                    |      UPDATE stock SET cantidad + X           |
                    |                                              |
                    |  FASE 2 - COMMIT o ROLLBACK                  |
                    |  Si ambas OK -> commit() en AMBOS nodos      |
                    |  Si alguna falla -> rollback() en AMBOS      |
                    +-----------------------------------------------+
\end{lstlisting}

\newpage
% ============================================================
\section{Procedimientos Almacenados}
% ============================================================

Los tres stored procedures están definidos en los archivos \code{sql/suc1.sql}, \code{suc2.sql} y \code{suc3.sql} (uno por nodo) y se crean automáticamente al inicializar los contenedores Docker.

\subsection{\texttt{sp\_realizar\_compra}}

\textbf{Parámetros:} \code{id\_sucursal}, \code{id\_cliente}, \code{id\_producto}, \code{cantidad}, \code{OUT resultado}, \code{OUT mensaje}

\textbf{Lógica:}
\begin{enumerate}[nosep]
    \item \code{SELECT cantidad ... FOR UPDATE} --- bloqueo de fila para evitar condiciones de carrera en compras simultáneas
    \item Valida que \code{cantidad\_disponible >= cantidad\_solicitada}; si no $\rightarrow$ \code{ROLLBACK}
    \item \code{INSERT INTO ventas} $\rightarrow$ obtiene \code{LAST\_INSERT\_ID()}
    \item \code{INSERT INTO detalle\_ventas}
    \item \code{UPDATE stock SET cantidad = cantidad - p\_cantidad}
    \item \code{COMMIT}
    \item \code{EXIT HANDLER FOR SQLEXCEPTION} $\rightarrow$ \code{ROLLBACK} automático ante cualquier error inesperado
\end{enumerate}

\textbf{Llamada desde PHP (\code{procesar\_venta.php}):}
\begin{lstlisting}[language=SQL]
CALL sp_realizar_compra(:suc, :cli, :prod, :cant, @resultado, @mensaje);
SELECT @resultado, @mensaje;
\end{lstlisting}

\subsection{\texttt{sp\_actualizar\_stock}}

\textbf{Parámetros:} \code{id\_sucursal}, \code{id\_producto}, \code{cantidad}, \code{operacion ('sumar'|'restar')}, \code{OUT resultado}, \code{OUT mensaje}

\textbf{Lógica:}
\begin{itemize}[nosep]
    \item Si no existe registro de stock $\rightarrow$ \code{INSERT} con cantidad inicial
    \item Si \code{operacion = 'restar'} y \code{stock < cantidad} $\rightarrow$ \code{ROLLBACK}
    \item Si \code{operacion = 'sumar'} $\rightarrow$ \code{UPDATE ... SET cantidad + p\_cantidad}
    \item Si \code{operacion = 'restar'} y hay stock suficiente $\rightarrow$ \code{UPDATE ... SET cantidad - p\_cantidad}
\end{itemize}

Usado para ajustes administrativos de inventario (por ejemplo, tras un conteo físico). No abre transacción a nivel de PHP: como el procedimiento ya maneja su propia transacción, \code{ajustar\_stock.php} lo invoca directamente sin envolverlo en un \code{beginTransaction()} adicional.

\textbf{Llamada desde PHP (\code{ajustar\_stock.php}):}
\begin{lstlisting}[language=SQL]
CALL sp_actualizar_stock(:suc, :prod, :cant, :op, @resultado, @mensaje);
SELECT @resultado, @mensaje;
\end{lstlisting}

Accesible desde el panel de administración, pestaña \textbf{``Ajuste de Stock''}.

\subsection{\texttt{sp\_reconstruir\_stock}}

\textbf{Parámetros:} \code{id\_sucursal}, \code{OUT resultado}, \code{OUT mensaje}

\textbf{Lógica:} Recalcula el stock de cada producto en una sucursal desde cero, aplicando la fórmula:

\[
\text{stock\_real} = \sum(\text{compras recibidas}) - \sum(\text{ventas completadas})
\]

\begin{lstlisting}[language=SQL]
UPDATE stock s
SET s.cantidad = (
    COALESCE(SELECT SUM(dc.cantidad) FROM detalle_compras ... WHERE estado='recibida', 0)
    -
    COALESCE(SELECT SUM(dv.cantidad) FROM detalle_ventas  ... WHERE estado='completada', 0)
)
WHERE s.id_sucursal = p_id_sucursal;
\end{lstlisting}

Se llama al recuperar un nodo tras una falla, para corregir posibles desincronizaciones ocurridas mientras estuvo OFFLINE.

\textbf{Llamada desde PHP (\code{reconstruir\_stock.php}):}
\begin{lstlisting}[language=SQL]
CALL sp_reconstruir_stock(:suc, @resultado, @mensaje);
SELECT @resultado, @mensaje;
\end{lstlisting}

\newpage
% ============================================================
\section{Aplicación del Teorema CAP}
% ============================================================

\subsection{Elección: CP (Consistencia + Tolerancia a Particiones)}

Ante una partición de red (un nodo deja de responder), el sistema prioriza la \textbf{consistencia} del inventario sobre la \textbf{disponibilidad total}.

\textbf{Justificación:} En un sistema de control de inventario físico, una inconsistencia (ej. stock descontado en un nodo pero no en otro) tiene costo operativo real: sobreventas, descuadres contables, pérdida de confianza. Una indisponibilidad temporal y acotada es preferible.

\subsection{Comportamiento según escenario}

\begin{table}[h]
\centering
\begin{tabularx}{\textwidth}{@{}X X l@{}}
\toprule
\textbf{Escenario} & \textbf{Comportamiento} & \textbf{Modelo CAP} \\
\midrule
3 nodos disponibles & CRUD, ventas y traspasos normales & --- \\
1 nodo caído, sin operaciones sobre él & Sistema opera; stock del nodo caído muestra ``N/D'' & Lectura disponible parcialmente \\
Venta en sucursal con nodo caído & Rechazada (HTTP 503), sin cambios en ningún nodo & \textbf{CP}: consistencia $>$ disponibilidad \\
Traspaso con origen o destino caído & Rollback 2PC completo, ningún nodo modificado & \textbf{CP}: integridad garantizada \\
Compra simultánea del mismo producto & \code{FOR UPDATE} serializa el acceso; no se duplica el descuento & \textbf{CP}: concurrencia controlada \\
Nodo simulado OFFLINE & \code{ConexionNodos::get()} lanza excepción antes de conectar & Falla controlada visible \\
\bottomrule
\end{tabularx}
\end{table}

\subsection{Comparación con otras estrategias}

\textbf{AP (Disponibilidad + Partición):} Permitiría completar ventas en nodos sanos aunque otros estén caídos, sincronizando después. Descartado porque durante la partición no hay forma de garantizar unicidad del stock (riesgo de sobreventa).

\textbf{CA (Consistencia + Disponibilidad):} Válido en red perfectamente confiable. No aplica a esta arquitectura LAN donde se simulan fallas reales.

\newpage
% ============================================================
\section{Manejo de Fallos}
% ============================================================

\subsection{Detección de Nodo Caído}

\code{ConexionNodos::get(id)} implementa dos niveles de detección:

\begin{enumerate}[nosep]
    \item \textbf{Falla simulada:} Lee \code{/tmp/libre\_mercado\_nodos.json}; si el nodo está marcado \code{offline}, lanza \code{RuntimeException} inmediatamente (sin intentar conexión TCP).
    \item \textbf{Falla real:} Si la conexión TCP al contenedor falla dentro de 1 segundo (\code{PDO::ATTR\_TIMEOUT = 1}), lanza \code{PDOException}.
\end{enumerate}

En ambos casos el código llamador captura \code{Throwable} y devuelve HTTP 503 con mensaje de error claro.

\subsection{Respuesta ante Falla por Tipo de Operación}

\textbf{Venta (\code{sp\_realizar\_compra}):}
\begin{lstlisting}
Nodo caido -> ConexionNodos::get() lanza excepcion
           -> procesar_venta.php captura -> HTTP 503
           -> El carrito no se vacia; el cliente puede reintentar en otra sucursal
           -> Stock intacto (no hubo ningun UPDATE)
\end{lstlisting}

\textbf{Traspaso 2PC:}
\begin{lstlisting}
Nodo origen OK, nodo destino CAIDO:
  Fase 1 -> beginTransaction() en origen OK
          -> beginTransaction() en destino -> excepcion
  Fase 2 -> rollBack() en origen (deshace el UPDATE de resta)
          -> HTTP 503 al cliente
          -> Ambos nodos quedan en estado anterior
\end{lstlisting}

\textbf{CRUD de catálogo (\code{MultiNodo}):}
\begin{lstlisting}
Cualquier nodo falla durante la replicacion:
  -> rollBack() en todos los nodos que llegaron a abrir transaccion
  -> HTTP 503 / mensaje de error
  -> El catalogo queda sin cambios en todos los nodos
\end{lstlisting}

\subsection{Simulación de Falla (Requisito 4)}

Desde el panel ``Simular Falla de Nodos'' en el panel de administración:

\begin{enumerate}[nosep]
    \item El administrador hace clic en \textbf{``Simular Falla''} sobre una sucursal
    \item Se hace \code{POST /api/nodo\_estado.php} con \code{\{id\_sucursal: X, estado: "offline"\}}
    \item El servidor escribe \code{\{"1":"online","2":"offline","3":"online"\}} en \code{/tmp/libre\_mercado\_nodos.json}
    \item A partir de ese momento, toda operación que requiera el nodo 2 recibe error controlado
\end{enumerate}

\textbf{Recuperación:}
\begin{enumerate}[nosep]
    \item Clic en \textbf{``Recuperar Nodo''} $\rightarrow$ estado vuelve a \code{online}
    \item Clic en \textbf{``Reconstruir Stock (sp)''} $\rightarrow$ llama a \code{sp\_reconstruir\_stock()} para recalcular inventario
\end{enumerate}

\subsection{Pruebas Obligatorias}

\begin{table}[h]
\centering
\begin{tabularx}{\textwidth}{@{}X X X@{}}
\toprule
\textbf{Prueba} & \textbf{Cómo reproducirla} & \textbf{Resultado esperado} \\
\midrule
Compra normal & Tienda $\rightarrow$ seleccionar producto $\rightarrow$ comprar & Venta exitosa, stock decrementado \\
Nodo apagado & Admin $\rightarrow$ ``Simular Falla'' en Suc. Centro $\rightarrow$ intentar comprar ahí & HTTP 503, mensaje claro, stock intacto \\
Recuperación nodo & Admin $\rightarrow$ ``Recuperar Nodo'' $\rightarrow$ ``Reconstruir Stock'' & Stock recalculado por \code{sp\_reconstruir\_stock} \\
Compra simultánea & Dos compras al mismo tiempo del mismo producto & \code{FOR UPDATE} serializa; sin duplicar descuento \\
Pérdida conexión (traspaso) & Simular falla en nodo destino $\rightarrow$ ejecutar traspaso & Rollback 2PC, ningún nodo modificado \\
\bottomrule
\end{tabularx}
\end{table}

\newpage
% ============================================================
\section{Conclusiones}
% ============================================================

El sistema \textbf{Libre Mercado Distribuido} demuestra que es posible implementar un modelo de comercio electrónico con tolerancia real a fallos de nodos en una red LAN, manteniendo la integridad de los datos en todo momento.

\textbf{Decisiones de diseño validadas:}

\begin{itemize}[nosep]
    \item \textbf{Modelo CP} es el adecuado para inventario físico: los datos incorrectos tienen mayor costo que la indisponibilidad temporal.
    \item \textbf{Stored Procedures} centralizan la lógica transaccional en la base de datos, asegurando atomicidad con \code{EXIT HANDLER FOR SQLEXCEPTION} y \code{FOR UPDATE} para concurrencia.
    \item \textbf{Two-Phase Commit manual} (PHP + PDO) permite operaciones distribuidas entre dos nodos sin depender de un coordinador externo.
    \item \textbf{Replicación de catálogo} mediante \code{MultiNodo.php} garantiza que productos, clientes y proveedores sean siempre consistentes entre sucursales.
    \item \textbf{Falla rápida} (\code{ATTR\_TIMEOUT = 1s}, lectura de archivo de estado) evita que una sucursal caída bloquee al resto del sistema indefinidamente.
\end{itemize}

\textbf{Limitaciones conocidas:}

\begin{itemize}[nosep]
    \item El coordinador PHP es un punto único de falla: si cae el servidor web, todo el sistema deja de funcionar. En producción se requeriría un balanceador de carga con múltiples instancias PHP.
    \item La sincronización post-falla es manual (\code{sp\_reconstruir\_stock}). En un sistema AP real se usaría replicación asíncrona automática.
    \item El 2PC implementado no es tolerante a fallas del propio coordinador durante la fase commit (problema de los generales bizantinos en 2PC).
\end{itemize}

\end{document}
